% =================================================================================
% Definitionen f�r rtm / rtp
% =================================================================================
\ifrtm
	\newcommand{\SADAProf}{Prof. Dr.-Ing. Ulrich Konigorski}
  \newcommand{\SADAinstitut}{Institut f�r Automatisierungstechnik und Mechatronik\\
                             Fachgebiet Regelungstechnik und Mechatronik\\
                             \SADAProf}
  \newcommand{\SADAwebsite}{www.rtm.tu-darmstadt.de}
  \newcommand{\SADAtel}{06151/16-4167}
  \newcommand{\SADAlogo}{\includegraphics[width=3cm]{./common/IAT_rtm}}
\else
  \newcommand{\SADAProf}{Prof. Dr.-Ing. Dr. h.c. Rolf Isermann}
  \newcommand{\SADAinstitut}{Institut f�r Automatisierungstechnik\\
                             und Mechatronik\\
                             FG Regelungstechnik und Prozessautomatisierung\\
                             \SADAProf}
  \newcommand{\SADAwebsite}{http://www.rtm.tu-darmstadt.de/rtp.html}
  \newcommand{\SADAtel}{06151/16-5465}
  \newcommand{\SADAlogo}{\includegraphics[scale=0.5]{./rtp_fuer_tex/rtp_fuer_tex_RGB.pdf}}	
\fi

% =================================================================================		
% Definitionen Vector
	\newcommand{\SADAlogoVector}{\includegraphics[scale=0.5]{./Bilder/vector-logo.pdf}}
	\newcommand{\SADAwebsiteVector}{https://vector.com/}
		
% =================================================================================
% Definitionen f�r die Erkl�rung, dass die Arbeit ohne Hilfe Dritter etc. erstellt wurde
% =================================================================================
% Die folgende Zeile deklariert die m�glichen Varianten der Erkl�rung, dass die
% Arbeit ohne Hilfe Dritter etc. erstellt wurde.
\def\ETIT{ETIT}\def\MBSA{MBSA}\def\MBDA{MBDA}
% =================================================================================


% =================================================================================
% Definitionen aus tudreport-Vorlage
% =================================================================================
\newif\ifTUDmargin
%\TUDmargintrue		% Sehr breiter rechter Rand
\TUDmarginfalse		% Schmaler (normaler) rechter Rand

\ifTUDmargin	% ggf. breiten Rand setzen
  \geometry{marginparsep=4.2mm,marginparwidth=28.64mm}
\fi

\newlength{\longtablewidth}
\setlength{\longtablewidth}{0.7\linewidth}
\addtolength{\longtablewidth}{-\marginparsep}
\addtolength{\longtablewidth}{-\marginparwidth}
% =================================================================================


% =================================================================================
% Farbige Deckseite, graue Balken auf allen anderen Seiten
% =================================================================================
\newif\ifOnlyColorFront
\OnlyColorFronttrue	    % Balken nur auf Deckblatt farbig
% \OnlyColorFrontfalse		% alle Balken farbig
% =================================================================================


% =================================================================================
% Befehle, die in scrreprt nicht exisiteren werden hier definiert, wenn diese
% Klasse verwendet wird.
% Auch die Seitenr�nder werden angepasst, so dass es grob wie mit tudreport
% aussieht.
% =================================================================================
\ifTUDdesign
	\title{\SADATitel}
	\subtitle{\SADAAutor}
	\subsubtitle{\SADATyp\ %-- \SADAAbgabe  
			\\Betreuer: \SADABetreuer{}, \SADABetreuerII{}
			\\Unternehmen: \SADAUnternehmen
			\\Abteilung: PES1}
	%\subsubsubtitle{\SADAUnternehmen}
	\setinstitutionlogo[height]{Bilder/vector_Logo_Typ_1}
	\geometry{left=30mm, right=20mm, top=15mm, bottom=20mm}
\else
	\title{\SADATitel}
	\subtitle{}
	\author{\SADAAutor}	% scrreprt erwartet Autor
	
	\newcommand{\subsubtitle}[1]{}
	\newcommand{\settitlepicture}[1]{}
	\newcommand{\printpicturesize}{}
	\newcommand{\institution}[1]{}
    \pagestyle{headings}

  % Diese Pakete nur extra einbinden, wenn NICHT tudreport als Basis.
	\usepackage{amssymb}
	\usepackage{geometry}
	
% \settitlepicture{}	% Bild f�r Deckblatt
% \printpicturesize
\fi
% =================================================================================


% =================================================================================
% Anpassung Absatzformat
% =================================================================================
% Der Abstand der Zeilen betr�gt das 1,25-fache des Standard-Abstands von
% LaTeX. Da technische Arbeiten viele Formeln und Bilder enthalten, werden
% Abs�tze durch einen zus�tzlichen vertikalen Zwischenraum statt durch einen
% Einzug getrennt.
\linespread{1.25}
\setlength{\parindent}{0mm}
\setlength{\parskip}{1ex}
% =================================================================================


% =================================================================================
% Texte f�r die R�ckseite des Titelblatts vorgeben
% =================================================================================
\uppertitleback{}
\lowertitleback{}
\dedication{}
% =================================================================================


% =================================================================================
% Informationen (Meta-Daten) f�r pdf
% =================================================================================
\hypersetup{
	pdftitle = {\SADATitel},
	pdfsubject = {},
	pdfauthor = {\SADAAutor},
	pdfkeywords = {},
	pdfcreator = {},
	pdfproducer = {LaTeX with hyperref},
	pdfstartview = {Fit},
	pdfpagelayout = {SinglePage}
}
% =================================================================================
